The Formalism of State-Transition Matching: A Mathematical Thesis on Proposal and Engagement Dynamics

Formal Logical Analysis Framework  
September 2026

This thesis supplies a complete, rigorous mathematical formulation of matching theory focused on the discrete states of Proposal and Engagement. It is grounded in cooperative game theory and discrete mathematics via the Gale–Shapley stable-marriage paradigm. The framework underpins algorithmic economics, network resource allocation, and decentralized multi-agent systems. It analyzes deferred acceptance, structural partial orders, and state-space convergence.

Contents
Introduction and Abstract  
Mathematical Definitions and Formal Foundations  
The Mechanics of Proposal and Engagement States  
Logical Ordering, Lattice Structures, and Proof of Convergence  

1. Introduction and Abstract
The work formalizes the terminology, logical ordering, and structural dynamics that govern Proposal and Engagement states. These concepts function as fundamental mechanisms in stable matching algorithms. The thesis examines the properties of deferred acceptance, structural partial orders, and convergence of the state space to a stable matching.

2. Mathematical Definitions and Formal Foundations

Let \( P = \{p_1, p_2, \dots, p_n\} \) be a finite set of \( n \) Proposers and \( R = \{r_1, r_2, \dots, r_m\} \) a finite set of \( m \) Reviewers.

Definition 2.1 (Preference Relation).  
For each agent \( a \in P \cup R \) there exists a strict, complete, and transitive linear order \( \succ_a \) over the opposing set of agents.  
Example: \( r_i \succ_{p_k} r_j \) denotes that proposer \( p_k \) strictly prefers reviewer \( r_i \) to reviewer \( r_j \).

Definition 2.2 (Matching Function).  
A matching \( \mu \) is a bijective (or one-to-one) mapping  
\[
\mu : P \cup R \to P \cup R \cup \{\emptyset\}
\]  
satisfying:
\( \mu(p) \in R \cup \{\emptyset\} \) for all \( p \in P \),  
\( \mu(r) \in P \cup \{\emptyset\} \) for all \( r \in R \),  
\( \mu(p) = r \iff \mu(r) = p \).

The empty set \( \emptyset \) represents an unmatched (null) assignment.

3. The Mechanics of Proposal and Engagement States

System dynamics evolve over discrete algorithmic time steps \( t \in \mathbb{N} \), tracking intentional proposals and conditional engagements.

Definition 3.1 (Proposal State Vector).  
A proposal \( \mathcal{P}_t(p, r) \) is a directed, intentional logical state emitted at time \( t \) by an unassigned agent \( p \in P \) and directed toward an agent \( r \in R \), where \( r \) is the maximal element under the restricted preference set of previously unproposed reviewers:
\[
r = \max_{\succ_p} \bigl\{ r' \in R \mid \forall t'  t_1 \), then \( p_2 \succeq_r p_1 \).

Theorem 4.2 (Finiteness and Convergence).  
The sequence of proposal–engagement sets terminates at a finite time step  
\[
T \le |P| \times |R|.
\]  
At step \( T \) the configuration stabilizes:  
\[
\mathcal{E}_T(p, r) \equiv \mu(p) = r,
\]  
yielding a Pareto-optimal stable matching.

Proof.  
Define the cumulative proposal volume
\[
V(t) = \sum_{p \in P} \bigl| \{ r \in R \mid \mathcal{P}_t(p, r) = 1 \} \bigr|.
\]  
Because \( P \) and \( R \) are finite, \( V(t) \) is strictly bounded above by \( |P| \times |R| \).  
In every time step \( t \) at which a stable matching has not yet been reached, at least one unengaged proposer must emit a new proposal, increasing \( V(t) \) by at least 1.  
Since \( V(t) \) is strictly monotonically increasing and bounded from above, the state sequence must reach a fixed point and terminate after finitely many steps.  
□

The terminal matching is stable (no blocking pair exists) and is the proposer-optimal stable matching produced by the classic deferred-acceptance procedure.

Structural Properties and Latent Axioms (Synthesis)

Proposer Downward Monotonicity** – successive proposals of any proposer move strictly downward in that proposer’s preference list.  
Reviewer Upward Monotonicity** – each successive engagement of a reviewer is at least as preferred as the previous one (Lemma 4.1).  
Deferred Acceptance** – engagements remain provisional until the global fixed point is reached, preventing premature commitment to locally suboptimal pairs.  
Bounded Iteration Space** – the entire history of proposals is confined to a finite set of cardinality at most \( |P| \times |R| \).  
Fixed-Point Termination** – the process converges to a Nash-stable matching in which no pair of agents can mutually improve by breaking their current assignment.

This formalism converts the informal notions of “making a proposal” and “entering an engagement” into precise, programmable state transitions of a discrete dynamical system, guaranteeing termination at a stable, Pareto-optimal matching.

Thesis Framework: The Mathematical Terminology of Project Proposals

This framework treats a project proposal not as a narrative business document but as a rigorous mathematical object: a constrained multi-dimensional optimization problem. Qualitative elements (scope, deliverables, timelines, risks, budgets) are mapped into formal notation so that evaluation becomes programmatic, verifiable, and free of semantic ambiguity. The goal is to maximize target value while remaining inside the intersection of the client’s demand domain and the provider’s feasible resource set.

1. Executive Abstract / Core Claim
Project proposals are modeled as mathematical entities rather than arbitrary text. By translating business objectives, resource variables, and operational criteria into structured formal notation, the framework:
Eliminates programmatic (and linguistic) ambiguity.
Optimizes capital and resource allocation.
Establishes mathematically verifiable boundaries for execution.
Converts subjective review into automated or semi-automated validation.

2. Foundation Axioms of Proposal Modeling
Every proposal is a multi-dimensional optimization function constrained by finite operational spaces. Three foundational pillars define the structure:

The Objective Function**: Maximize target value yield \( f(x) \) relative to strategic milestones.  
  \( f(x) \) encodes the value the client seeks (revenue, capability, risk reduction, etc.).

The Resource Manifold**: A bounded vector space representing real-world limitations—capital, human bandwidth, temporal capacity, and other scarce inputs. Denote this manifold \( M \) (or the resource matrix/space \( R \)).

The Constraint Mapping**: Systemic equations (equalities and inequalities) that enforce structural, financial, regulatory, and operational compliance. These define the feasible region.

A proposal is admissible if and only if its solution lies inside the intersection of the client’s demand domain and the provider’s resource capability space.

3. Core Mathematical Lexicon (Translation Table)
Proposal components map directly onto formal terms to remove semantic drift:

| Proposal Element   | Mathematical Equivalent          | Operational Definition |
|--------------------|----------------------------------|------------------------|
| Project Scope      | Vector Domain \( (S) \)         | Strict boundary containing all acceptable operational states. |
| Deliverables       | Discrete Set Members \( (\mathbb{D}) \) | Finite elements that must be generated to satisfy the objective function. |
| Project Timeline   | Linear Continuum / Directed Bounded Interval \( (T) \) | Bounded directed interval mapping execution progression over time. |
| Risk Profile       | Stochastic Multiplier / Stochastic Variance Coefficient | Probability distributions reflecting variance in projected yields. |
| Budget Allocation  | Resource Matrix \( (R) \) or Boundary Constraints \( g(x) \le c \) | Transformation tensor (or linear inequalities) distributing capital across vector dimensions and preventing over-allocation. |

Additional mappings appearing in the operational view:
Resource bottlenecks → identified via Lagrangian multipliers.
Shadow prices → marginal cost of relaxing any single constraint.

4. Operational Mechanics & Optimization
Evaluation is treated as a Pareto Optimization Problem.  

A proposal is admissible if and only if it lies at the intersection of the client’s demand domain and the provider’s resource capability space (i.e., on or inside the feasible Pareto set).  
Resource bottlenecks are isolated using the Lagrangian Multiplier system.  
  Form the Lagrangian  
  \[
  \mathcal{L}(x,\lambda) = f(x) + \sum_i \lambda_i g_i(x)
  \]
  (or the appropriate inequality form). The multipliers \( \lambda \) (shadow prices) quantify the precise marginal cost of expanding any given resource constraint (staff capacity, deadline, capital ceiling, etc.).  

This converts proposal drafting and review from subjective judgment into precise, programmable optimization: one can compute how much additional value would be obtained by relaxing a particular constraint and decide whether that expansion is justified.

5. Structural Conclusion
Moving from linguistic frameworks to formal mathematical notation:
Standardizes proposal structures.
Reduces human error and interpretive variance.
Enables automated computational review and programmatic validation.
Turns subjective evaluation into an objective, auditable process.

The resulting framework serves as a foundational architectural blueprint for technical writers, procurement officers, quantitative business analysts, and anyone who must draft, evaluate, or negotiate proposals under resource constraints.

Practical Implications for “Engaging in a Proposal”
When writing or assessing a proposal under this paradigm:
Explicitly define the objective function \( f(x) \) and the vector domain \( S \).
Enumerate deliverables as the discrete set \( \mathbb{D} \).
Map the timeline to the bounded interval \( T \).
Encode risks as stochastic multipliers and budgets as the resource matrix / inequality constraints \( g(x) \le c \).
Solve (or approximate) the constrained optimization problem; examine the Pareto frontier and the shadow prices given by the Lagrangian multipliers.
Accept the proposal only if a feasible point exists inside the intersection of demand and capability; otherwise identify the binding constraints and their associated shadow prices as negotiation levers.

This is the complete, self-contained mathematical terminology and operational mechanics of the thesis.

Thesis Framework: SELORT + Proper-LLM-Diction-Output  
Logical Skeleton & Verification Order Coupled with Controlled Expressive Surface

This dual-layer paradigm governs the generation, verification, and rendering of claims, proposals, arguments, or any assertoric text. It cleanly separates what must be proven and in what irreversible sequence from how the verified content is linguistically realized. Nothing is asserted until it has traversed a strict syntactic-entailment pipeline; only then is the proven content rendered under quantifiable stylistic control that simultaneously maximizes richness, variation, and anti-repetition.

The two complementary systems are:
SELORT** (Syntactic Entailment Logical Order Realization Theorem) — the logical skeleton and verification order (https://github.com/TheAstrographer/SELORT).
Proper-LLM-Diction-Output** — the controlled expressive surface (https://github.com/TheAstrographer/Proper-LLM-Diction-Output).

Together they convert subjective or probabilistic generation into a two-stage process of (1) irreversible logical verification followed by (2) steerable, high-diction expression. Earlier single-layer frameworks (Pareto/resource-manifold optimization or Gale–Shapley state-transition matching) become optional rather than mandatory.

1. Executive Abstract
SELORT supplies the irreversible logical skeleton. Proper-LLM-Diction-Output supplies the measurable stylistic surface.  

The combined system enforces:
Proof-first ordering**: a claim may be treated as true only after it has traversed the canonical chain.
Controlled linguistic realization**: once verified, the content is rendered with explicit stylistic parameters that prevent flat or repetitive prose while remaining faithful to the logical core.

2. Foundational Pillars

Pillar A – Logical Skeleton (SELORT)  
Every claim is treated as a set of formulas \(\Gamma\) that must realize entailment before any further step is permitted. The canonical chain is irreversible and left-to-right:

\[
\ulcorner\Gamma \rightsquigarrow \vdash \rightsquigarrow \models_p \rightsquigarrow T\urcorner
\]

or, in multi-witness form:

\[
\ulcorner\vdash \varphi_1,\varphi_2,\dots,\varphi_n = \psi \rightsquigarrow \models_p \psi \rightsquigarrow T\psi\urcorner
\]

\(\Gamma\) itself realizes entailment in the first place.  
Realization (\(\rightsquigarrow\)) is strictly downstream from \(\Gamma\).  
Model-theoretic shortcuts (\(\models_s\)) are blocked.  
Truth (\(T\)) is the terminal confirmation that the entailment has become semantic.

Pillar B – Expressive Surface (Proper-LLM-Diction-Output)  
Once a statement has been confirmed true under the SELORT chain, it is passed to a stylistic controller defined by the six-dimensional functional vector

\[
\overline{V} = [C,\, D,\, R,\, M,\, \Delta,\, \rho]
\]

where:
\(C\) = Cadence (staccato vs. sinuous balance),
\(D\) = Diction density (lexical uniqueness / type-token ratio),
\(R\) = Rhythm (normalized average sentence length),
\(M\) = Metaphor / figurative load,
\(\Delta\) = Register deviation (variance in complexity and length),
\(\rho\) = Surmounted (elevated / complex) diction density.

This vector, together with dynamic stylistic operators, density-aware word banks, and negative-space triggers, governs how the verified logical content is rendered.

3. Core Operational Mechanics

Stage 1 – Verification Order (SELORT / ProofFirstTruthEngine)  
For any candidate claim \(\psi\) relative to premises \(\Gamma\):

gamma_realizes(Γ, ψ) — \(\Gamma\) itself must realize the entailment.  
prove(Γ, ψ) — Syntactic proof judgment \(\vdash\) is recorded.  
realize(Γ, ψ) — Internal realization map is activated (blocked unless prior steps succeeded).  
entail_p(Γ, ψ) — Proof-entailment judgment \(\models_p\).  
confirm_truth(Γ, ψ) — Terminal truth confirmation \(T\).

Any attempt to short-circuit the order raises an explicit block. Multi-formula witness sets are supported. Even deterministic computational sequences (e.g., seeded PRNG streams) are treated as derived syntactic objects that must still traverse the same chain.

Stage 2 – Expressive Rendering (Proper-LLM-Diction-Output)  
Only after Stage 1 returns true is the content eligible for surface realization. The rendering engine:
Computes or targets a desired \(\overline{V}\).
Selects vocabulary from density-aware word banks.
Applies cadence shaping (staccato fragmentation or sinuous flow).
Injects controlled metaphor load and register variation.
Optionally inserts negative-space (typographic or structural pauses).
Can be further constrained by an edit-distance (e.g., Levenshtein) budget relative to a canonical formal phrasing, ensuring the surface does not drift beyond a measurable threshold from the verified logical core.

4. Integrated Lexicon & Translation Table

| Conceptual Element     | SELORT Term / Mechanism                  | Diction-Output Term / Mechanism       | Combined Role                                      |
|------------------------|------------------------------------------|---------------------------------------|----------------------------------------------------|
| Premises / Ground      | \(\Gamma\) (non-empty set of formulas)  | Base lexical material                 | Starting point that must realize entailment        |
| Claim to be verified   | \(\psi\)                                 | Core propositional content            | Object of both proof and rendering                 |
| Proof step             | \(\vdash\)                               | —                                     | Syntactic acceptance gate                          |
| Realization            | \(\rightsquigarrow\)                     | Style-vector activation               | Transition from logic to expression                |
| Proof-entailment       | \(\models_p\)                            | —                                     | Semantic elevation of the proof                    |
| Terminal truth         | \(T\)                                    | Verified content ready for \(\overline{V}\) | Only point at which rich rendering is permitted |
| Surface variation      | —                                        | \(C, D, R, M, \Delta, \rho\)          | Controlled richness & anti-repetition              |
| Faithfulness constraint| Irreversible chain                       | Optional Levenshtein / similarity budget | Prevents stylistic drift from logical core      |

5. Structural Advantages & Operational Consequences
Separation of concerns**: Logic is never contaminated by stylistic flourish; style is never applied to unverified claims.  
Irreversibility**: Once a step in the SELORT chain is passed, earlier steps cannot be revisited or weakened.  
Measurability**: Both the logical status (boolean progression through the chain) and the linguistic quality (\(\overline{V}\) values) are explicit and programmable.  
Anti-hallucination posture**: Nothing is asserted until \(\Gamma\) has realized the entailment and the full chain has completed.  
Anti-repetition posture**: The style vector and word-bank mechanisms actively vary surface form while remaining anchored to the verified content.  
Extensibility**: Formal proof assistants (Lean 4, Coq, Isabelle/HOL, Metamath, Agda) can be attached via the proof-first diction adapters; any downstream rendering engine can consume the confirmed \(T\) statements.

6. Structural Conclusion
By coupling SELORT’s strict, proof-first verification order with Proper-LLM-Diction-Output’s quantifiable expressive surface, the framework achieves a complete pipeline:

Prove first, under irreversible syntactic order → then render, under controlled stylistic richness.

Claims, proposals, or any assertoric text can be processed solely through this SELORT + Diction pipeline. The earlier single-layer approaches (Mathematical Terminology of Project Proposals with Pareto frontiers / Lagrangian multipliers, or the Formalism of State-Transition Matching) are thereby rendered optional. The result is text that is both logically certified and expressively optimized—a foundational blueprint for technical writing, formal argumentation, proposal generation, and any domain that demands simultaneous rigor and linguistic vitality.

Final statement of the thesis  
SELORT supplies the logical skeleton and the ordered sequence of what must be proven before anything may be asserted; Proper-LLM-Diction-Output supplies the expressive surface that renders the proven content with controlled richness, variation, and anti-repetition.
